\begin{solapartat}
\punts{0,2} Nombre màssic = A = 99

\smallskip
\punts{0,2} Nombre atòmic = Protons = Z = 43

\smallskip
\punts{0,2} Neutrons = N = A -- Z = 56

\smallskip
\punts{0,4} ${}^{99}_{43}Tc^{*} \rightarrow {}^{99}_{43}Tc + \gamma$
\end{solapartat}

\begin{solapartat}
\punts{0,3} $m = m_0 e^{-kt}$

\smallskip
\punts{0,3} $\dfrac{m_0}{2} = m_0 e^{-kt_{1/2}} \rightarrow \ln\!\left(\dfrac{1}{2}\right) = -kt_{1/2} \rightarrow k = \dfrac{\ln 2}{6} = 0,116\,h^{-1}$

\smallskip
\punts{0,4} $m = 2\cdot\underbrace{e^{-\left(\frac{\ln 2}{6}\right)\cdot 24}}_{0,0625} = 0,125\,ng$

\medskip
També es pot fer argumentant que per cada període de semidesintegració la mostra queda reduïda a la meitat. Com que passen 4 períodes la mostra queda reduïda a la setzena part. $m = \frac{m_0}{16} = 0,125\,ng$
\end{solapartat}
